% --- Archon physics formalization source begin ---
% archon:physics
% archon:covers PhyXMiniProblems/problem_phyx_mini_0236.lean
% archon:source-report reports/phyx_mini/problem_phyx_mini_0236.source.json

\chapter{Physics problem phyx\_mini\_0236}
\label{ch:PhyXMiniProblems_problem_phyx_mini_0236}

\paragraph{Problem source.}
\#\# Physical scenario

An object of mass \textbackslash{}( m\_1 = 9.00\textbackslash{},\textbackslash{}mathrm\{kg\} \textbackslash{}) is in equilibrium when connected to a light spring of constant \textbackslash{}( k = 100\textbackslash{},\textbackslash{}mathrm\{N/m\} \textbackslash{}) that is fastened to a wall as shown in figure. A second object, \textbackslash{}( m\_2 = 7.00\textbackslash{},\textbackslash{}mathrm\{kg\} \textbackslash{}), is slowly pushed up against \textbackslash{}( m\_1 \textbackslash{}), compressing the spring by the amount \textbackslash{}( A = 0.200\textbackslash{},\textbackslash{}mathrm\{m\} \textbackslash{}) (see figure). The system is then released, and both objects start moving to the right on the frictionless surface.

\#\# Question

How far apart are the objects when the spring is fully stretched for the first time?

\#\# Answer choices

- A: A: 8.86cm

- B: B: 8.76cm

- C: C: 8.66cm

- D: D: 8.56cm

\#\# Figure caption (auxiliary; use the image as primary evidence)

The image consists of four panels labeled (a) through (d), showing a mechanical system with a spring and two blocks:

1. **Panel (a):**
   - A spring with a spring constant \textbackslash{}( k \textbackslash{}) is attached to a wall on the left side.
   - The spring is connected to a block labeled \textbackslash{}( m\_1 \textbackslash{}), which is stationary at an equilibrium position.

2. **Panel (b):**
   - The spring is compressed, and block \textbackslash{}( m\_1 \textbackslash{}) is moved to the left.
   - A second block \textbackslash{}( m\_2 \textbackslash{}), is adjacent to \textbackslash{}( m\_1 \textbackslash{}), and both blocks are positioned at a distance \textbackslash{}( A \textbackslash{}) to the left of the equilibrium position.

3. **Panel (c):**
   - Upon release of the spring, both blocks \textbackslash{}( m\_1 \textbackslash{}) and \textbackslash{}( m\_2 \textbackslash{}) move to the right.
   - There's a red arrow labeled \textbackslash{}( \textbackslash{}vec\{v\} \textbackslash{}) indicating the direction of velocity.

4. **Panel (d):**
   - Block \textbackslash{}( m\_1 \textbackslash{}) remains attached to the spring, which is now stretched.
   - Block \textbackslash{}( m\_2 \textbackslash{}) continues moving to the right alone, as shown by the red arrow with velocity \textbackslash{}( \textbackslash{}vec\{v\} \textbackslash{}).
   - The distance \textbackslash{}( D \textbackslash{}) represents the separation between \textbackslash{}( m\_1 \textbackslash{}) and \textbackslash{}( m\_2 \textbackslash{}).

These panels illustrate the process of compressing and releasing a spring, causing the blocks to move sequentially.

\#\# Dataset metadata

Wave Properties; Multi-Formula Reasoning, Physical Model Grounding Reasoning

\paragraph{Recorded answer/context.}
D: D: 8.56cm

\paragraph{Figure/image path.}
/root/proposal\_for\_physic/science-mango/phyx\_mini\_run/phyx\_data/test\_image/236.png

\paragraph{Formalization target.}
create a compiling Lean file with sorry bodies at `PhyXMiniProblems/problem\_phyx\_mini\_0236.lean`. The Lean declarations must preserve the physical quantities, dimensions or dimensional roles, figure labels, governing-law hypotheses, and final relation expressed by this problem.
Use Mathlib/Physlib names found through LeanExplore where available. If a physics API is missing, introduce faithful local abstractions rather than scalar placeholder aliases.

\begin{theorem}[Physics formalization target]
\label{thm:physics:phyx_mini_0236:target}
\lean{PhyXMiniProblems.ProblemPhyXMini0236.firstFullyStretchedSeparation_matches_answerD}
\uses{def:physics:phyx-mini-0236:phyxminiproblems-problemphyxmini0236-lengthinmeters, def:physics:phyx-mini-0236:phyxminiproblems-problemphyxmini0236-twoblockspringsetup, def:physics:phyx-mini-0236:phyxminiproblems-problemphyxmini0236-matchesproblemandprimaryfigure, def:physics:phyx-mini-0236:phyxminiproblems-problemphyxmini0236-hasphysicalparameters, def:physics:phyx-mini-0236:phyxminiproblems-problemphyxmini0236-isfirstmaximumspringextension, def:physics:phyx-mini-0236:phyxminiproblems-problemphyxmini0236-satisfiestwostagespringdynamics, def:physics:phyx-mini-0236:phyxminiproblems-problemphyxmini0236-matchesanswerchoice}
The assigned autoformalize agent should translate this physics problem into Lean declarations in the covered file, with theorem and lemma proof bodies written as `by sorry`.
\end{theorem}
\begin{proof}
This is an autoformalization task, not a proof task. Produce statements that can later be proved by the physics prover without weakening the physical model.
\end{proof}

% --- Archon declaration topology begin ---
\section{Declaration topology}
\begin{definition}[Mass Quantity]
  \label{def:physics:phyx-mini-0236:phyxminiproblems-problemphyxmini0236-massquantity}
  \lean{PhyXMiniProblems.ProblemPhyXMini0236.MassQuantity}
  A physical mass
\end{definition}
\begin{proof}
  This is the declared model definition of Mass Quantity.
\end{proof}
\begin{definition}[Length Quantity]
  \label{def:physics:phyx-mini-0236:phyxminiproblems-problemphyxmini0236-lengthquantity}
  \lean{PhyXMiniProblems.ProblemPhyXMini0236.LengthQuantity}
  A physical length or signed one-dimensional displacement
\end{definition}
\begin{proof}
  This is the declared model definition of Length Quantity.
\end{proof}
\begin{definition}[Duration Quantity]
  \label{def:physics:phyx-mini-0236:phyxminiproblems-problemphyxmini0236-durationquantity}
  \lean{PhyXMiniProblems.ProblemPhyXMini0236.DurationQuantity}
  A physical duration
\end{definition}
\begin{proof}
  This is the declared model definition of Duration Quantity.
\end{proof}
\begin{definition}[Speed Quantity]
  \label{def:physics:phyx-mini-0236:phyxminiproblems-problemphyxmini0236-speedquantity}
  \lean{PhyXMiniProblems.ProblemPhyXMini0236.SpeedQuantity}
  A one-dimensional speed magnitude, with dimension length per time
\end{definition}
\begin{proof}
  This is the declared model definition of Speed Quantity.
\end{proof}
\begin{definition}[Spring Constant Quantity]
  \label{def:physics:phyx-mini-0236:phyxminiproblems-problemphyxmini0236-springconstantquantity}
  \lean{PhyXMiniProblems.ProblemPhyXMini0236.SpringConstantQuantity}
  A spring constant has dimension force per length, equivalently mass per time squared. Its SI readout is measured in newtons per metre
\end{definition}
\begin{proof}
  This is the declared model definition of Spring Constant Quantity.
\end{proof}
\begin{definition}[mass Readout]
  \label{def:physics:phyx-mini-0236:phyxminiproblems-problemphyxmini0236-massreadout}
  \lean{PhyXMiniProblems.ProblemPhyXMini0236.massReadout}
  \uses{def:physics:phyx-mini-0236:phyxminiproblems-problemphyxmini0236-massquantity}
  Read a mass in a selected mass unit
\end{definition}
\begin{proof}
  This is the declared model definition of mass Readout.
\end{proof}
\begin{definition}[length Readout]
  \label{def:physics:phyx-mini-0236:phyxminiproblems-problemphyxmini0236-lengthreadout}
  \lean{PhyXMiniProblems.ProblemPhyXMini0236.lengthReadout}
  \uses{def:physics:phyx-mini-0236:phyxminiproblems-problemphyxmini0236-lengthquantity}
  Read a length or displacement in a selected length unit
\end{definition}
\begin{proof}
  This is the declared model definition of length Readout.
\end{proof}
\begin{definition}[duration Readout]
  \label{def:physics:phyx-mini-0236:phyxminiproblems-problemphyxmini0236-durationreadout}
  \lean{PhyXMiniProblems.ProblemPhyXMini0236.durationReadout}
  \uses{def:physics:phyx-mini-0236:phyxminiproblems-problemphyxmini0236-durationquantity}
  Read a duration in a selected time unit
\end{definition}
\begin{proof}
  This is the declared model definition of duration Readout.
\end{proof}
\begin{definition}[speed Readout]
  \label{def:physics:phyx-mini-0236:phyxminiproblems-problemphyxmini0236-speedreadout}
  \lean{PhyXMiniProblems.ProblemPhyXMini0236.speedReadout}
  \uses{def:physics:phyx-mini-0236:phyxminiproblems-problemphyxmini0236-speedquantity}
  Read a speed in coherent selected length and time units
\end{definition}
\begin{proof}
  This is the declared model definition of speed Readout.
\end{proof}
\begin{definition}[spring Constant Readout]
  \label{def:physics:phyx-mini-0236:phyxminiproblems-problemphyxmini0236-springconstantreadout}
  \lean{PhyXMiniProblems.ProblemPhyXMini0236.springConstantReadout}
  \uses{def:physics:phyx-mini-0236:phyxminiproblems-problemphyxmini0236-springconstantquantity}
  Read a spring constant in coherent selected mass and time units
\end{definition}
\begin{proof}
  This is the declared model definition of spring Constant Readout.
\end{proof}
\begin{definition}[mass In Kilograms]
  \label{def:physics:phyx-mini-0236:phyxminiproblems-problemphyxmini0236-massinkilograms}
  \lean{PhyXMiniProblems.ProblemPhyXMini0236.massInKilograms}
  \uses{def:physics:phyx-mini-0236:phyxminiproblems-problemphyxmini0236-massquantity, def:physics:phyx-mini-0236:phyxminiproblems-problemphyxmini0236-massreadout}
  Kilogram readout used for the two stated masses
\end{definition}
\begin{proof}
  This is the declared model definition of mass In Kilograms.
\end{proof}
\begin{definition}[length In Meters]
  \label{def:physics:phyx-mini-0236:phyxminiproblems-problemphyxmini0236-lengthinmeters}
  \lean{PhyXMiniProblems.ProblemPhyXMini0236.lengthInMeters}
  \uses{def:physics:phyx-mini-0236:phyxminiproblems-problemphyxmini0236-lengthquantity, def:physics:phyx-mini-0236:phyxminiproblems-problemphyxmini0236-lengthreadout}
  Metre readout used for the compression and the governing laws
\end{definition}
\begin{proof}
  This is the declared model definition of length In Meters.
\end{proof}
\begin{definition}[length In Centimeters]
  \label{def:physics:phyx-mini-0236:phyxminiproblems-problemphyxmini0236-lengthincentimeters}
  \lean{PhyXMiniProblems.ProblemPhyXMini0236.lengthInCentimeters}
  \uses{def:physics:phyx-mini-0236:phyxminiproblems-problemphyxmini0236-lengthquantity, def:physics:phyx-mini-0236:phyxminiproblems-problemphyxmini0236-lengthreadout}
  Centimetre readout used by the displayed answers
\end{definition}
\begin{proof}
  This is the declared model definition of length In Centimeters.
\end{proof}
\begin{definition}[duration In Seconds]
  \label{def:physics:phyx-mini-0236:phyxminiproblems-problemphyxmini0236-durationinseconds}
  \lean{PhyXMiniProblems.ProblemPhyXMini0236.durationInSeconds}
  \uses{def:physics:phyx-mini-0236:phyxminiproblems-problemphyxmini0236-durationquantity, def:physics:phyx-mini-0236:phyxminiproblems-problemphyxmini0236-durationreadout}
  Second readout used for elapsed time after the blocks cease touching
\end{definition}
\begin{proof}
  This is the declared model definition of duration In Seconds.
\end{proof}
\begin{definition}[speed In Meters Per Second]
  \label{def:physics:phyx-mini-0236:phyxminiproblems-problemphyxmini0236-speedinmeterspersecond}
  \lean{PhyXMiniProblems.ProblemPhyXMini0236.speedInMetersPerSecond}
  \uses{def:physics:phyx-mini-0236:phyxminiproblems-problemphyxmini0236-speedquantity, def:physics:phyx-mini-0236:phyxminiproblems-problemphyxmini0236-speedreadout}
  Metres-per-second readout of a speed magnitude
\end{definition}
\begin{proof}
  This is the declared model definition of speed In Meters Per Second.
\end{proof}
\begin{definition}[spring Constant In Newtons Per Meter]
  \label{def:physics:phyx-mini-0236:phyxminiproblems-problemphyxmini0236-springconstantinnewtonspermeter}
  \lean{PhyXMiniProblems.ProblemPhyXMini0236.springConstantInNewtonsPerMeter}
  \uses{def:physics:phyx-mini-0236:phyxminiproblems-problemphyxmini0236-springconstantquantity, def:physics:phyx-mini-0236:phyxminiproblems-problemphyxmini0236-springconstantreadout}
  Newton-per-metre readout of the spring constant
\end{definition}
\begin{proof}
  This is the declared model definition of spring Constant In Newtons Per Meter.
\end{proof}
\begin{definition}[Block]
  \label{def:physics:phyx-mini-0236:phyxminiproblems-problemphyxmini0236-block}
  \lean{PhyXMiniProblems.ProblemPhyXMini0236.Block}
  The two blocks labeled in panels (a)--(d)
\end{definition}
\begin{proof}
  This is the declared model definition of Block.
\end{proof}
\begin{definition}[Figure Panel]
  \label{def:physics:phyx-mini-0236:phyxminiproblems-problemphyxmini0236-figurepanel}
  \lean{PhyXMiniProblems.ProblemPhyXMini0236.FigurePanel}
  The four panels in the supplied image
\end{definition}
\begin{proof}
  This is the declared model definition of Figure Panel.
\end{proof}
\begin{definition}[Horizontal Direction]
  \label{def:physics:phyx-mini-0236:phyxminiproblems-problemphyxmini0236-horizontaldirection}
  \lean{PhyXMiniProblems.ProblemPhyXMini0236.HorizontalDirection}
  Horizontal arrow directions in the figure
\end{definition}
\begin{proof}
  This is the declared model definition of Horizontal Direction.
\end{proof}
\begin{definition}[Spring State]
  \label{def:physics:phyx-mini-0236:phyxminiproblems-problemphyxmini0236-springstate}
  \lean{PhyXMiniProblems.ProblemPhyXMini0236.SpringState}
  Qualitative state of the spring relative to its equilibrium length
\end{definition}
\begin{proof}
  This is the declared model definition of Spring State.
\end{proof}
\begin{definition}[Figure Distance Mark]
  \label{def:physics:phyx-mini-0236:phyxminiproblems-problemphyxmini0236-figuredistancemark}
  \lean{PhyXMiniProblems.ProblemPhyXMini0236.FigureDistanceMark}
  Distance labels printed under panels (b) and (d)
\end{definition}
\begin{proof}
  This is the declared model definition of Figure Distance Mark.
\end{proof}
\begin{definition}[Surface Condition]
  \label{def:physics:phyx-mini-0236:phyxminiproblems-problemphyxmini0236-surfacecondition}
  \lean{PhyXMiniProblems.ProblemPhyXMini0236.SurfaceCondition}
  Surface idealizations relevant to the horizontal motion
\end{definition}
\begin{proof}
  This is the declared model definition of Surface Condition.
\end{proof}
\begin{definition}[Spring Idealization]
  \label{def:physics:phyx-mini-0236:phyxminiproblems-problemphyxmini0236-springidealization}
  \lean{PhyXMiniProblems.ProblemPhyXMini0236.SpringIdealization}
  Idealization of the spring stated in the problem
\end{definition}
\begin{proof}
  This is the declared model definition of Spring Idealization.
\end{proof}
\begin{definition}[Two Block Spring Figure]
  \label{def:physics:phyx-mini-0236:phyxminiproblems-problemphyxmini0236-twoblockspringfigure}
  \lean{PhyXMiniProblems.ProblemPhyXMini0236.TwoBlockSpringFigure}
  \uses{def:physics:phyx-mini-0236:phyxminiproblems-problemphyxmini0236-block, def:physics:phyx-mini-0236:phyxminiproblems-problemphyxmini0236-figurepanel, def:physics:phyx-mini-0236:phyxminiproblems-problemphyxmini0236-horizontaldirection, def:physics:phyx-mini-0236:phyxminiproblems-problemphyxmini0236-springstate, def:physics:phyx-mini-0236:phyxminiproblems-problemphyxmini0236-figuredistancemark}
  Qualitative evidence transcribed from the four-panel primary image. The boolean distance marks record only the appearances of A and D; their physical values live in the setup and are not inferred from pixels
\end{definition}
\begin{proof}
  This is the declared model definition of Two Block Spring Figure.
\end{proof}
\begin{definition}[Two Block Spring Setup]
  \label{def:physics:phyx-mini-0236:phyxminiproblems-problemphyxmini0236-twoblockspringsetup}
  \lean{PhyXMiniProblems.ProblemPhyXMini0236.TwoBlockSpringSetup}
  \uses{def:physics:phyx-mini-0236:phyxminiproblems-problemphyxmini0236-massquantity, def:physics:phyx-mini-0236:phyxminiproblems-problemphyxmini0236-lengthquantity, def:physics:phyx-mini-0236:phyxminiproblems-problemphyxmini0236-durationquantity, def:physics:phyx-mini-0236:phyxminiproblems-problemphyxmini0236-speedquantity, def:physics:phyx-mini-0236:phyxminiproblems-problemphyxmini0236-springconstantquantity, def:physics:phyx-mini-0236:phyxminiproblems-problemphyxmini0236-block, def:physics:phyx-mini-0236:phyxminiproblems-problemphyxmini0236-horizontaldirection, def:physics:phyx-mini-0236:phyxminiproblems-problemphyxmini0236-surfacecondition, def:physics:phyx-mini-0236:phyxminiproblems-problemphyxmini0236-springidealization, def:physics:phyx-mini-0236:phyxminiproblems-problemphyxmini0236-twoblockspringfigure}
  Independent physical quantities for the two stages of the motion. The displacement functions use the spring equilibrium position as origin. After contact ends, m1DisplacementAfterContactEnd t is the rightward signed displacement of m₁, while m2DisplacementAfterContactEnd t is the distance traveled freely by m₂. The separation is stored independently and is related to those displacements only by the geometry law below
\end{definition}
\begin{proof}
  This is the declared model definition of Two Block Spring Setup.
\end{proof}
\begin{definition}[Matches Problem And Primary Figure]
  \label{def:physics:phyx-mini-0236:phyxminiproblems-problemphyxmini0236-matchesproblemandprimaryfigure}
  \lean{PhyXMiniProblems.ProblemPhyXMini0236.MatchesProblemAndPrimaryFigure}
  \uses{def:physics:phyx-mini-0236:phyxminiproblems-problemphyxmini0236-massinkilograms, def:physics:phyx-mini-0236:phyxminiproblems-problemphyxmini0236-lengthinmeters, def:physics:phyx-mini-0236:phyxminiproblems-problemphyxmini0236-speedinmeterspersecond, def:physics:phyx-mini-0236:phyxminiproblems-problemphyxmini0236-springconstantinnewtonspermeter, def:physics:phyx-mini-0236:phyxminiproblems-problemphyxmini0236-twoblockspringsetup}
  Numerical data from the prose and qualitative facts from the primary image. Panel (a) supplies the equilibrium reference; (b) shows compression A and touching blocks; (c) shows their common rightward motion at equilibrium; and (d) shows stretched m₁, freely moving m₂, and the unknown distance D. No numerical value for D occurs here
\end{definition}
\begin{proof}
  This is the declared model definition of Matches Problem And Primary Figure.
\end{proof}
\begin{definition}[Has Physical Parameters]
  \label{def:physics:phyx-mini-0236:phyxminiproblems-problemphyxmini0236-hasphysicalparameters}
  \lean{PhyXMiniProblems.ProblemPhyXMini0236.HasPhysicalParameters}
  \uses{def:physics:phyx-mini-0236:phyxminiproblems-problemphyxmini0236-massinkilograms, def:physics:phyx-mini-0236:phyxminiproblems-problemphyxmini0236-lengthinmeters, def:physics:phyx-mini-0236:phyxminiproblems-problemphyxmini0236-durationinseconds, def:physics:phyx-mini-0236:phyxminiproblems-problemphyxmini0236-speedinmeterspersecond, def:physics:phyx-mini-0236:phyxminiproblems-problemphyxmini0236-springconstantinnewtonspermeter, def:physics:phyx-mini-0236:phyxminiproblems-problemphyxmini0236-twoblockspringsetup}
  Positivity and nondegeneracy conditions for the physical experiment
\end{definition}
\begin{proof}
  This is the declared model definition of Has Physical Parameters.
\end{proof}
\begin{definition}[Is First Maximum Spring Extension]
  \label{def:physics:phyx-mini-0236:phyxminiproblems-problemphyxmini0236-isfirstmaximumspringextension}
  \lean{PhyXMiniProblems.ProblemPhyXMini0236.IsFirstMaximumSpringExtension}
  \uses{def:physics:phyx-mini-0236:phyxminiproblems-problemphyxmini0236-durationquantity, def:physics:phyx-mini-0236:phyxminiproblems-problemphyxmini0236-lengthinmeters, def:physics:phyx-mini-0236:phyxminiproblems-problemphyxmini0236-durationinseconds, def:physics:phyx-mini-0236:phyxminiproblems-problemphyxmini0236-twoblockspringsetup}
  The selected time is the first time at which m₁ reaches a global maximum rightward extension after the blocks cease touching. Global maximality says the spring is fully stretched; strict inequality at every earlier nonnegative time selects its first occurrence
\end{definition}
\begin{proof}
  This is the declared model definition of Is First Maximum Spring Extension.
\end{proof}
\begin{definition}[Satisfies Two Stage Spring Dynamics]
  \label{def:physics:phyx-mini-0236:phyxminiproblems-problemphyxmini0236-satisfiestwostagespringdynamics}
  \lean{PhyXMiniProblems.ProblemPhyXMini0236.SatisfiesTwoStageSpringDynamics}
  \uses{def:physics:phyx-mini-0236:phyxminiproblems-problemphyxmini0236-durationquantity, def:physics:phyx-mini-0236:phyxminiproblems-problemphyxmini0236-massreadout, def:physics:phyx-mini-0236:phyxminiproblems-problemphyxmini0236-lengthreadout, def:physics:phyx-mini-0236:phyxminiproblems-problemphyxmini0236-speedreadout, def:physics:phyx-mini-0236:phyxminiproblems-problemphyxmini0236-springconstantreadout, def:physics:phyx-mini-0236:phyxminiproblems-problemphyxmini0236-massinkilograms, def:physics:phyx-mini-0236:phyxminiproblems-problemphyxmini0236-lengthinmeters, def:physics:phyx-mini-0236:phyxminiproblems-problemphyxmini0236-durationinseconds, def:physics:phyx-mini-0236:phyxminiproblems-problemphyxmini0236-speedinmeterspersecond, def:physics:phyx-mini-0236:phyxminiproblems-problemphyxmini0236-springconstantinnewtonspermeter, def:physics:phyx-mini-0236:phyxminiproblems-problemphyxmini0236-twoblockspringsetup}
  Governing laws for the two motion stages. Mechanical energy is conserved while the touching blocks move together from the compressed release state to the equilibrium/contact-end state. Physlib's harmonic oscillator is tied to the SI readouts of m₁ and k. After contact ends, m₁ follows the SHO solution with initial velocity v, while m₂ moves uniformly at the same speed. Their separation is the free-block displacement minus the spring-block displacement.
\end{definition}
\begin{proof}
  This is the declared model definition of Satisfies Two Stage Spring Dynamics.
\end{proof}
\begin{lemma}[contact End Speed eq compression mul sqrt stiffness div total Mass]
  \label{lem:physics:phyx-mini-0236:phyxminiproblems-problemphyxmini0236-contactendspeed-eq-compression-mul-sqrt-stiffness-div-totalmass}
  \lean{PhyXMiniProblems.ProblemPhyXMini0236.contactEndSpeed_eq_compression_mul_sqrt_stiffness_div_totalMass}
  \uses{def:physics:phyx-mini-0236:phyxminiproblems-problemphyxmini0236-massinkilograms, def:physics:phyx-mini-0236:phyxminiproblems-problemphyxmini0236-lengthinmeters, def:physics:phyx-mini-0236:phyxminiproblems-problemphyxmini0236-speedinmeterspersecond, def:physics:phyx-mini-0236:phyxminiproblems-problemphyxmini0236-springconstantinnewtonspermeter, def:physics:phyx-mini-0236:phyxminiproblems-problemphyxmini0236-twoblockspringsetup, def:physics:phyx-mini-0236:phyxminiproblems-problemphyxmini0236-matchesproblemandprimaryfigure, def:physics:phyx-mini-0236:phyxminiproblems-problemphyxmini0236-hasphysicalparameters, def:physics:phyx-mini-0236:phyxminiproblems-problemphyxmini0236-satisfiestwostagespringdynamics}
  Energy conservation during the common-motion stage gives the speed at the instant contact ends and the spring reaches equilibrium
\end{lemma}
\begin{proof}
  The conclusion follows from the governing relations and figure readouts named in the statement of contact End Speed eq compression mul sqrt stiffness div total Mass.
\end{proof}
\begin{lemma}[first Maximum Stretch Time eq quarter Period]
  \label{lem:physics:phyx-mini-0236:phyxminiproblems-problemphyxmini0236-firstmaximumstretchtime-eq-quarterperiod}
  \lean{PhyXMiniProblems.ProblemPhyXMini0236.firstMaximumStretchTime_eq_quarterPeriod}
  \uses{def:physics:phyx-mini-0236:phyxminiproblems-problemphyxmini0236-durationinseconds, def:physics:phyx-mini-0236:phyxminiproblems-problemphyxmini0236-twoblockspringsetup, def:physics:phyx-mini-0236:phyxminiproblems-problemphyxmini0236-hasphysicalparameters, def:physics:phyx-mini-0236:phyxminiproblems-problemphyxmini0236-isfirstmaximumspringextension, def:physics:phyx-mini-0236:phyxminiproblems-problemphyxmini0236-satisfiestwostagespringdynamics}
  The first global maximum of the post-contact sine trajectory occurs one quarter of an m₁ oscillator period after contact ends
\end{lemma}
\begin{proof}
  The conclusion follows from the governing relations and figure readouts named in the statement of first Maximum Stretch Time eq quarter Period.
\end{proof}
\begin{lemma}[first Maximum Separation formula]
  \label{lem:physics:phyx-mini-0236:phyxminiproblems-problemphyxmini0236-firstmaximumseparation-formula}
  \lean{PhyXMiniProblems.ProblemPhyXMini0236.firstMaximumSeparation_formula}
  \uses{def:physics:phyx-mini-0236:phyxminiproblems-problemphyxmini0236-massinkilograms, def:physics:phyx-mini-0236:phyxminiproblems-problemphyxmini0236-lengthinmeters, def:physics:phyx-mini-0236:phyxminiproblems-problemphyxmini0236-twoblockspringsetup, def:physics:phyx-mini-0236:phyxminiproblems-problemphyxmini0236-matchesproblemandprimaryfigure, def:physics:phyx-mini-0236:phyxminiproblems-problemphyxmini0236-hasphysicalparameters, def:physics:phyx-mini-0236:phyxminiproblems-problemphyxmini0236-isfirstmaximumspringextension, def:physics:phyx-mini-0236:phyxminiproblems-problemphyxmini0236-satisfiestwostagespringdynamics}
  Combining the common contact-end speed, the quarter-period SHO motion, free motion of m₂, and one-dimensional separation geometry eliminates k and gives the exact general separation formula
\end{lemma}
\begin{proof}
  The conclusion follows from the governing relations and figure readouts named in the statement of first Maximum Separation formula.
\end{proof}
\begin{definition}[Answer Choice]
  \label{def:physics:phyx-mini-0236:phyxminiproblems-problemphyxmini0236-answerchoice}
  \lean{PhyXMiniProblems.ProblemPhyXMini0236.AnswerChoice}
  Labels of the four displayed answer choices
\end{definition}
\begin{proof}
  This is the declared model definition of Answer Choice.
\end{proof}
\begin{definition}[displayed Centimeters]
  \label{def:physics:phyx-mini-0236:phyxminiproblems-problemphyxmini0236-answerchoice-displayedcentimeters}
  \lean{PhyXMiniProblems.ProblemPhyXMini0236.AnswerChoice.displayedCentimeters}
  \uses{def:physics:phyx-mini-0236:phyxminiproblems-problemphyxmini0236-answerchoice}
  Separation in centimetres printed beside each answer label
\end{definition}
\begin{proof}
  This is the declared model definition of displayed Centimeters.
\end{proof}
\begin{definition}[Matches Answer Choice]
  \label{def:physics:phyx-mini-0236:phyxminiproblems-problemphyxmini0236-matchesanswerchoice}
  \lean{PhyXMiniProblems.ProblemPhyXMini0236.MatchesAnswerChoice}
  \uses{def:physics:phyx-mini-0236:phyxminiproblems-problemphyxmini0236-lengthquantity, def:physics:phyx-mini-0236:phyxminiproblems-problemphyxmini0236-lengthincentimeters, def:physics:phyx-mini-0236:phyxminiproblems-problemphyxmini0236-answerchoice}
  Agreement with a separation displayed to the nearest hundredth of a centimetre. This rounding predicate is answer-list metadata, not a premise
\end{definition}
\begin{proof}
  This is the declared model definition of Matches Answer Choice.
\end{proof}
% --- Archon declaration topology end ---
% --- Archon physics formalization source end ---
